\section{Color Management and the Working Space}
\label{sec:colorspace}

\subsection{RAW Decode Configuration}
\label{sec:raw-config}

The correctness of every subsequent stage rests on obtaining a genuinely
scene-referred linear image. Apple's RAW pipeline \cite{coreimage} defaults are
tuned to produce a pleasing display-referred result, and every one of those
defaults is destructive to our purposes. The required configuration is:

\begin{lstlisting}[language=Swift, caption={Scene-referred RAW decode configuration.}, label={lst:rawconfig}]
func makeSceneReferredFilter(url: URL) throws -> CIFilter & CIRAWFilter {
    guard let f = CIRAWFilter(imageURL: url) else { throw DecodeError.unsupported }

    // 1. Defeat every rendering intent Apple applies by default.
    f.isGamutMappingEnabled       = false
    f.localToneMapAmount          = 0.0
    f.extendedDynamicRangeAmount  = 0.0   // we manage headroom ourselves
    f.boostAmount                 = 0.0   // disables the S-curve entirely
    f.boostShadowAmount           = 0.0
    f.contrastAmount              = 0.0
    f.shadowBias                  = 0.0

    // 2. Defeat spatial processing we intend to model ourselves.
    f.sharpnessAmount             = 0.0
    f.detailAmount                = 0.0
    f.contrastAmount              = 0.0
    f.luminanceNoiseReductionAmount = 0.0
    f.colorNoiseReductionAmount     = 0.0

    // 3. Keep the vendor's geometric and shading corrections: these are
    //    calibration, not interpretation, and reimplementing them is
    //    value-destructive.
    f.isLensCorrectionEnabled     = true

    // 4. Preserve capture white balance as the *starting point*; the user
    //    may override, but we must know what the camera decided.
    f.neutralTemperature          = f.neutralTemperature
    f.neutralTint                 = f.neutralTint

    return f
}
\end{lstlisting}

The decoded image is requested in a linear working color space with an
extended range, and the render destination is configured with
\texttt{CIContext} options that disable any output color management, because
the render graph performs its own conversion:

\begin{lstlisting}[language=Swift, caption={Context configuration for linear extended-range output.}]
let ws = CGColorSpace(name: CGColorSpace.extendedLinearDisplayP3)!
let ctx = CIContext(mtlDevice: device, options: [
    .workingColorSpace  : ws,
    .workingFormat      : CIFormat.RGBAf,
    .cacheIntermediates : false,
    .highQualityDownsample : true
])
\end{lstlisting}

\subsection{Choice of Working Space}

The pipeline requires a working space that (i) encloses the gamut of every
iPhone sensor after color matrix application, (ii) encloses the gamut of the
modelled film dyes, which extend outside Rec.\,709 in the cyan and yellow
regions, (iii) has a white point compatible with the printing model, and (iv)
has a documented, invertible relationship to CIE XYZ \cite{hunt2011}.

Three candidates were evaluated:

\begin{table}[htbp]
\caption{Working Space Evaluation}
\label{tab:workingspace}
\begin{center}
\renewcommand{\arraystretch}{1.15}
\begin{tabularx}{\columnwidth}{@{}l>{\raggedright\arraybackslash}X@{}}
\toprule
\textbf{Space} & \textbf{Assessment} \\
\midrule
Linear Display P3 & Native to the platform, zero-cost preview conversion, but clips saturated cyans that the modelled cyan dye reaches. Rejected as primary; retained as the display target. \\
ACEScg (AP1) & Wide enough, industry standard, well-documented transforms, D60 white point requires adaptation to D65 for display. Chosen. \\
DaVinci Wide Gamut & Wider still, excellent for interchange with Resolve, but the transform is vendor-published rather than standardized. Retained as an export option only. \\
\bottomrule
\end{tabularx}
\end{center}
\end{table}

\EM{} adopts \textbf{ACEScg} \cite{acescg,st2065} (ACES AP1 primaries, linear,
D60 white) as the
internal scene-referred working space. The choice costs one matrix multiply on
input and one on output relative to staying in Display P3, and buys gamut
headroom that prevents negative-valued channels from appearing mid-pipeline ---
which matters enormously, because $\logE$ of a negative number is undefined and
a single clamped channel produces a visible hue shift in saturated highlights.

\subsection{Input Transform}

The decoded image arrives in extended linear Display P3. The conversion to
ACEScg is the composition of the P3-to-XYZ matrix, a Bradford
\cite{lam1985} chromatic
adaptation from D65 to D60, and the XYZ-to-AP1 matrix:

\begin{equation}
\mathbf{M}_{\mathrm{in}} = \mathbf{M}_{\mathrm{XYZ}\to\mathrm{AP1}}\,
\mathbf{M}_{\mathrm{CAT}}^{\mathrm{D65}\to\mathrm{D60}}\,
\mathbf{M}_{\mathrm{P3}\to\mathrm{XYZ}}
\label{eq:min}
\end{equation}

evaluating numerically to

\begin{equation}
\mathbf{M}_{\mathrm{in}} \approx
\begin{pmatrix}
\phantom{-}0.9525 & \phantom{-}0.0343 & \phantom{-}0.0132 \\
\phantom{-}0.0170 & \phantom{-}0.9754 & \phantom{-}0.0076 \\
-0.0018 & \phantom{-}0.0107 & \phantom{-}0.9911
\end{pmatrix}.
\label{eq:minval}
\end{equation}

The matrix is precomputed on the host at profile-load time, not in the shader.

The inverse transform for the display path is $\mathbf{M}_{\mathrm{out}} =
\mathbf{M}_{\mathrm{in}}^{-1}$ followed by the Display P3 transfer function.

\subsection{Exposure Normalization and the Anchor Problem}

Film responds to absolute exposure; a digital sensor reports a relative value
scaled by an unknown gain. Before the characteristic curve can be applied, the
image must be anchored so that a known scene reflectance maps to a known
position on the curve.

\EM{} adopts the standard cinematographic anchor: an 18\% diffuse reflector
under correct exposure maps to a relative exposure of $E = 0.18$ in the working
space, which is placed at the film's rated speed point. The transformation is

\begin{equation}
E_{\mathrm{anchored}} = E_{\mathrm{decoded}} \cdot 2^{\,\varepsilon} \cdot
\frac{0.18}{g_{\mathrm{cal}}}
\label{eq:anchor}
\end{equation}

where $\varepsilon$ is the user's exposure compensation in stops and
$g_{\mathrm{cal}}$ is a per-device calibration constant relating the decoded
middle-gray value to the working-space unit. Determining $g_{\mathrm{cal}}$ is a
calibration task, specified in Section~\ref{sec:calibration}; a wrong
$g_{\mathrm{cal}}$ manifests as every image sitting too high or too low on the
characteristic curve, which the user experiences as ``the stocks all look
flat'' or ``everything blocks up.''

The relation between anchored exposure and the film's rated ISO speed $S$ is

\begin{equation}
\logE_{\mathrm{film}} = \log_{10}\!\left(E_{\mathrm{anchored}}\right)
+ \log_{10}\!\left(\frac{S}{S_0}\right) + x_{\mathrm{ref}}
\label{eq:isoshift}
\end{equation}

with $S_0 = 100$ the reference speed and $x_{\mathrm{ref}}$ the stock's
speed-point log exposure. This is what makes ``shooting Portra 400 at 800''
meaningful: it shifts $\logE_{\mathrm{film}}$ by $-0.301$, moving the entire
image down the characteristic curve into the toe, which the user then
compensates for by pushing development (Section~\ref{sec:development}) ---
exactly the real relationship, produced by the model rather than scripted.

\subsection{White Balance in the Working Space}

White balance is applied as a von Kries adaptation in the LMS cone space rather
than as an RGB channel gain, because channel gain in a wide-gamut space
produces hue rotation in saturated colors. The operator is

\begin{equation}
\mathbf{E}' = \mathbf{M}_{\mathrm{AP1}\to\mathrm{LMS}}^{-1}\,
\mathrm{diag}\!\left(\frac{\boldsymbol{\rho}_{\mathrm{dst}}}{\boldsymbol{\rho}_{\mathrm{src}}}\right)
\mathbf{M}_{\mathrm{AP1}\to\mathrm{LMS}}\, \mathbf{E}
\label{eq:vonkries}
\end{equation}

where $\boldsymbol{\rho}$ are the cone responses to the source and destination
illuminants and $\mathbf{M}_{\mathrm{AP1}\to\mathrm{LMS}}$ uses the CAT02
\cite{cie159}
transform. Because both matrices and the diagonal are constant per frame, the
whole of \eqref{eq:vonkries} collapses to a single $3\times3$ matrix computed on
the host, so the shader cost is one matrix multiply.

\subsection{Tungsten Stocks and the Illuminant Question}

A tungsten-balanced stock such as Vision3 500T is not a color filter. It is an
emulsion whose three layers have different relative sensitivities, calibrated so
that a 3200\,K illuminant produces equal densities. Modelling it as a blue push
is wrong; the correct model is a per-layer speed offset $\Delta x_c$ applied to
the log exposure before the characteristic curve:

\begin{equation}
x_c \leftarrow x_c + \Delta x_c^{\mathrm{(balance)}}, \quad c \in \{R,G,B\}
\label{eq:tungsten}
\end{equation}

with $\Delta x^{\mathrm{(balance)}}$ derived from the ratio of the stock's
layer sensitivities integrated against the reference illuminant SPD. For a
tungsten stock shot under daylight without correction, this produces the
familiar strong blue cast --- and, importantly, produces it \emph{with the
correct nonlinearity}, because the blue record is driven far up its curve while
red sits low, so the cast varies with scene luminance rather than being a
uniform tint. Table~\ref{tab:balance} gives the shipping values.

\begin{table}[htbp]
\caption{Layer Balance Offsets $\Delta x^{\mathrm{(balance)}}$}
\label{tab:balance}
\begin{center}
\renewcommand{\arraystretch}{1.15}
\begin{tabular}{@{}lccc@{}}
\toprule
\textbf{Stock class} & $\Delta x_R$ & $\Delta x_G$ & $\Delta x_B$ \\
\midrule
Daylight (5500\,K aim)  & $0.000$ & $0.000$ & $0.000$ \\
Tungsten (3200\,K aim)  & $-0.310$ & $-0.055$ & $+0.365$ \\
Type L long-exposure    & $-0.285$ & $-0.040$ & $+0.325$ \\
Monochrome panchromatic & $0.000$ & $0.000$ & $0.000$ \\
\bottomrule
\end{tabular}
\end{center}
\end{table}
